% MoVi (gruppo SDMs): documento di due pagine per la giuria, versione 3.0.
% Composto da sciame-comunicazione (f20-chiusura/31-proposta) il 15-11-2026 sulla v2.0 (copia intatta: proposta_sciame_v2.0.tex).
% Cifre e frasi SOLO da f20-chiusura/20-verdetto/VERDETTO_FINALE.md (sigle V-, S-, LP-, CR- e tabelle 4 e 5 del verdetto);
% la tabella con la sigla di ogni cifra sta in f20-chiusura/31-proposta/PROPOSTA_V3.md.
% Si compila con: xelatex proposta_sciame.tex
\documentclass[9pt,a4paper]{extarticle}
\usepackage[a4paper,margin=1.2cm,top=0.8cm,bottom=0.8cm]{geometry}
\usepackage{fontspec}
\setmainfont{Fira Sans}[Scale=0.93,BoldFont={Fira Sans SemiBold},ItalicFont={Fira Sans Italic}]
\usepackage{polyglossia}\setmainlanguage{italian}
\usepackage{xcolor,tikz,tabularx,booktabs,array,enumitem,microtype,ragged2e,amsmath}
\usetikzlibrary{positioning,arrows.meta}
\usepackage[hidelinks]{hyperref}

\definecolor{ink}{HTML}{16202B}
\definecolor{muted}{HTML}{55606C}
\definecolor{accent}{HTML}{0B5C8A}
\definecolor{accentbg}{HTML}{E8F1F6}
\definecolor{warnbg}{HTML}{FBF1E3}
\definecolor{ok}{HTML}{1E7A4C}
\definecolor{warn}{HTML}{A8620A}
\definecolor{bad}{HTML}{A12B2B}
\definecolor{rule}{HTML}{C9D1D9}
\color{ink}

\newcommand{\versione}{v3.1}
\newcommand{\dataora}{15 novembre 2026}
\newcommand{\statodoc}{volo ad alta velocità; con documento tecnico}

% Classi del verdetto finale: lettera + colore, leggibili anche in bianco e nero
\newcommand{\chip}[2]{\tikz[baseline=(c.base)]\node[inner xsep=2.2pt,inner ysep=1.0pt,rounded corners=1.5pt,fill=#1,text=white,font=\fontsize{6.0}{6.0}\selectfont\bfseries](c){#2};}
\newcommand{\Sv}{\chip{accent}{S}}   % solido
\newcommand{\Iv}{\chip{ok}{I}}       % incoraggiante
\newcommand{\Lv}{\chip{warn}{L}}     % limite dichiarato
\newcommand{\Nv}{\chip{bad}{NP}}     % non provato

\setlength{\parindent}{0pt}\setlength{\parskip}{1.4pt}
\setlist{leftmargin=1.0em,itemsep=0.3pt,topsep=0.8pt,parsep=0pt,partopsep=0pt}
\newcommand{\sez}[1]{\vspace{1.8pt}{\color{accent}\bfseries\fontsize{9.4}{10.2}\selectfont #1}\par\vspace{-3.8pt}{\color{rule}\rule{\linewidth}{0.5pt}}\par\vspace{0pt}}
\renewcommand{\arraystretch}{1.08}
\newcolumntype{L}{>{\RaggedRight\arraybackslash}X}
\newcommand{\nota}[1]{{\color{muted}\fontsize{7.5}{8.9}\selectfont #1}}
\newcommand{\w}[1]{{\color{muted}[#1]}}
\pagestyle{empty}
\emergencystretch=1.5em

\begin{document}
\fontsize{8.5}{9.7}\selectfont

% ---------------------------------------------------------------- testata
\begin{minipage}[b]{0.79\linewidth}
{\fontsize{15.5}{17.5}\selectfont\bfseries\color{accent} MoVi: una percezione che dichiara ciò che non ha osservato}\\[2pt]
{\fontsize{9.0}{10.8}\selectfont \textbf{Aerei e droni in volo veloce, con collisioni, detriti ed eventi improvvisi.} Dove la vista si perde, MoVi non dice mai «libero»: dichiara dove l'oggetto può essere e quanto vale, secondo dopo secondo, quella regione.}\\[1pt]
{\fontsize{8.0}{9.8}\selectfont\color{muted}\textbf{Gruppo SDMs:} Ignazio Crocchiolo, Federico Dellini, Danilo Amalfitano, Emanuele Ciotola, Pierfrancesco Amendola · The Big Hack 2026}
\end{minipage}\hfill
\begin{minipage}[b]{0.19\linewidth}\raggedleft
{\bfseries\color{accent}\versione}\; · \;\dataora\\
\nota{\statodoc}
\end{minipage}

\vspace{2pt}{\color{accent}\rule{\linewidth}{1.1pt}}\vspace{0.3pt}

\nota{\Sv\ solido \; · \; \Iv\ incoraggiante: vicino al sicuro, con base matematica \; · \; \Lv\ limite dichiarato \; · \; \Nv\ non provato. Stato di ogni cifra: \textit{sim} simulato dalle schede; \textit{pub} su traiettorie pubbliche con sensore simulato; \textit{calc} calcolato; \textit{ip} ipotesi; \textit{dim} dimostrato nel modello. Le cifre sono passaggi simulati (k su n) con l'intervallo di Wilson al 95\,\% per passaggio, fra parentesi quadre.}

% ---------------------------------------------------------------- limiti
\vspace{0.5pt}
\colorbox{warnbg}{\parbox{\dimexpr\linewidth-2\fboxsep}{\vspace{0.3pt}\textbf{Limiti, da leggere per primi}
\begin{itemize}
\item \textbf{Tutto è simulato:} dalle schede tecniche, da un calcolo, o su traiettorie di volo pubbliche con un sensore simulato. Nessun sensore fisico, nessuna prova su hardware; niente è validato. Le classi degli oggetti sono messe a mano (il classificatore si suppone esistente). Ogni scenario ha una sola geometria nominale.
\item \textbf{La collisione veloce fra droni non è ancora dichiarata come evento} \Lv: il nucleo adottato la riconosce in 0 passaggi su 60 \w{0; 0,060}. MoVi la vede arrivare e tiene la regione (§\,2); una regola nuova è in prova, non adottata (§\,4).
\item \textbf{A bordo di un drone, con il radar accecato dall'esplosione} l'avviso al portatore crolla (0 su 59 nel caso peggiore) \Lv; una correzione in prova lo riporta a 59 su 59 \w{0,939; 1}, ma non localizza l'evento e dà 11 falsi avvisi su 200 \w{0,031; 0,096} se il radar satura per altra causa. L'innesco della fiammata è ottimista (saturazione presa dalla finestra vera dell'evento).
\item MoVi osserva, monitora e avvisa: nessuna uscita è un comando, nessuna funzione d'ingaggio, di puntamento o di contrasto.
\end{itemize}\vspace{0.3pt}}}

% ---------------------------------------------------------------- 1
\sez{1 · L'idea: tre risposte, e mai «libero» su ciò che non si è visto}
Per ogni oggetto MoVi dà una di tre risposte: \textbf{lo vedo}; \textbf{non lo vedo, ma è in questa regione}; \textbf{non lo so}. «Libero» vale solo dove si è osservato. L'avviso di possibile presenza si spegne solo guardando tutta la zona, non per il tempo. Quando un sensore è saturato da una fiammata o accecato, MoVi lo dichiara «non affidabile», e da quel sensore non viene più nessun «libero».

\begin{center}
\begin{tikzpicture}[
  b/.style={draw=accent,line width=0.55pt,rounded corners=2pt,fill=white,align=center,inner xsep=2pt,inner ysep=2pt,font=\fontsize{7.2}{8.4}\selectfont,minimum height=10mm,text width=27.0mm},
  f/.style={-{Stealth[length=3.6pt]},line width=0.55pt,draw=ink}, node distance=2.6mm]
\node[b](s){\textbf{Sensori} (da scheda)\\radar a scansione 10\,Hz · radar rotante · camere visibile e termica · ADS-B};
\node[b,right=of s](h){\textbf{Salute dei sensori}\\fiammata, fumo, cecità: «non affidabile»};
\node[b,right=of h](k){\textbf{Tracce e fusione}\\Kalman a più regimi (IMM), conferma su più cicli};
\node[b,right=of k,fill=accentbg](r){\textbf{Regione e confidenza}\\$C(\tau)$ che cala; insieme di presenza a $v_{\max}$};
\node[b,right=of r](e){\textbf{Eventi e detriti}\\avaria, collisione: dedotti; volume dei detriti};
\node[b,right=of e](u){\textbf{Uscita}\\vedo / regione / non so; registro «osservato» o «dedotto»};
\foreach \a/\bb in {s/h,h/k,k/r,r/e,e/u} \draw[f](\a)--(\bb);
\end{tikzpicture}
\end{center}
\vspace{-5pt}

% ---------------------------------------------------------------- 2
\sez{2 · Il caso centrale: volo ad alta velocità con eventi improvvisi}
Sei scenari simulati con sensori di classe attuale (radar a scansione elettronica e radar rotante): avaria con fiammata, due droni veloci in rotta di collisione, quasi collisioni, incroci veloci, ottici accecati, radar rotante. Ogni cifra è \textit{sim}, 60 passaggi per cella (200 per gli incroci).
\begin{itemize}
\item \Sv\ \textbf{Mai «libero» dove c'era l'oggetto o un frammento}, nei sei scenari: 0 su 200 \w{0; 0,019} negli incroci veloci, 0 su 60 \w{0; 0,060} nelle altre celle. Base: «libero» solo ciò che è osservato; insieme di presenza a 210 m/s.
\item \Iv\ \textbf{Avaria in volo di un aereo a 135 m/s con fiammata:} evento dedotto entro 1 s e camere dichiarate non affidabili in 60 su 60 \w{0,940; 1} (innesco ottimista). I detriti a bassa energia restano tutti nel volume dichiarato in 58 su 60 \w{0,886; 0,991}. L'avviso arriva prima del primo frammento a terra in 60 su 60 \w{0,940; 1}; il margine minimo per un operatore è di 0,25–0,45 s (risposta umana da FAA AC 90-48E: \textit{ip}).
\item \Iv\ \textbf{Due droni veloci in rotta di collisione:} avviso di prossimità almeno 10 s prima in 60 su 60 \w{0,940; 1} (soglia dei 10 s non scritta prima). \Sv\ \textbf{Dopo l'urto} il punto vero e i frammenti sono sempre rimasti nella regione dichiarata: 60 su 60 \w{0,940; 1}, 939 campioni su 939 (inviluppo con $|v|\le 210$ m/s).
\item \Sv\ \textbf{Col nucleo adottato, nessun falso evento} nelle quasi collisioni a 40 m e negli incroci veloci: 0 su 200 \w{0; 0,019}, 0 su 60. \Sv\ \textbf{Ottici accecati per 2–4 s:} ogni traccia conserva una regione, 60 su 60 \w{0,940; 1}; l'identità resta «presunta» \Lv\ (scambio 0,47–0,63).
\item \Sv\ \textbf{Radar a scansione elettronica a 10 Hz:} la regione non si perde mai, 60 su 60 e 200 su 200 \w{0,981; 1}. \Lv\ Col radar rotante a 1 Hz l'aereo veloce perde continuità, 39 su 60 \w{0,524; 0,758}, e MoVi lo dichiara.
\end{itemize}

% ---------------------------------------------------------------- 3
\sez{3 · Quando la continuità si perde: la confidenza nel tempo, in forma chiusa}
\begin{minipage}[t]{0.44\linewidth}
\vspace{0pt}
Dal momento della perdita, a velocità quasi costante con rumore gaussiano (\textit{dim} nel modello, verificata numericamente; un regime, nel piano: la copertura in tabella è misurata in 3D, con miscela e ricalibrazione):
\[\sigma^2(\tau)=\sigma_0^2+2\sigma_{0v}\tau+\sigma_v^2\tau^2+\tfrac{q\tau^3}{3}+\sigma_{\mathrm{rif}}^2\]
\[C(\tau)=1-e^{-r^2/2\sigma^2(\tau)},\qquad r_{0,9}(\tau)=2{,}146\,\sigma(\tau)\]
Con più regimi di moto $C=\sum_j\mu_j P_j$; si mostra il minimo progressivo $\tilde C(\tau)=\min_{s\le\tau}C(s)$, che non risale mai. Fuori modello (urto, avaria) si mostra l'insieme di presenza $r_0+v_{\max}\tau$ (210 m/s aerei, 45 m/s droni): un contenimento garantito se $|v|\le v_{\max}$, non una probabilità.
\end{minipage}\hfill
\begin{minipage}[t]{0.54\linewidth}
\vspace{1pt}
{\fontsize{7.8}{9.2}\selectfont
\begin{tabularx}{\linewidth}{@{}>{\RaggedRight}p{2.3cm}LLL@{}}
\toprule
Aerei, regione al 90\,\% (\textit{pub}) & 1–2 s & 21–30 s & in virata \\
\midrule
prima della taratura & 0,71 \w{0,67; 0,74} & 0,96 \w{0,94; 0,98} & 0,905 (stima) \\
dopo la taratura & \Iv\ \textbf{0,89} \w{0,87; 0,91} & \Iv\ 0,85 \w{0,81; 0,88} & \Lv\ 0,73 \w{0,67; 0,78} \\
$r_{0,9}$ mediano & 39 m & 1106 m & — \\
\bottomrule
\end{tabularx}}
\nota{Dopo la perdita la regione al 90\,\% ha contenuto l'aereo da 85 a 89 volte su 100 fra 1 e 30 s. I casi sono correlati (intervalli ottimisti) e vengono per nove decimi da quattro voli della stessa rotta. In virata la copertura scende a 73 su 100, in accelerazione a 0,76 \w{0,67; 0,83}: limiti dichiarati. Taratura su metà dei mezzi, prova sull'altra metà. Per il drone seguito di §\,4 (\textit{calc} da \textit{ip}): $C$ scende da 1,00 a 0,62 a 2 s e a 0,09 a 5 s (radar frontale).}
\end{minipage}

% ---------------------------------------------------------------- 4
\sez{4 · MoVi a bordo di un drone: la traccia mantenuta attraverso una collisione vicina}
Il sensore viaggia su un drone che mantiene la traccia di un drone veloce più avanti. A 30–100 m due droni si scontrano: fiammata, fumo, 32 frammenti. \textit{sim}, 60 passaggi per geometria (di lato, davanti, dietro, in virata), più 20 per cella su semi e geometrie nuovi.
\begin{itemize}
\item \Sv\ \textbf{La regione che contiene il drone seguito non si è mai persa:} 60 su 60 \w{0,940; 1} in ogni geometria, col radar a 360°; 20 su 20 \w{0,839; 1} anche con l'urto a 10–30 m (robustezza, 20 per cella). \textbf{Mai «libero»} sul drone seguito né su un frammento: 0 su 60 \w{0; 0,060} in ogni cella. Nessun frammento attribuito al drone seguito senza dirlo: l'identità diventa «presunta, dubbia».
\item \Sv\ \textbf{Il volume dei detriti dichiarato ha contenuto tutti i frammenti simulati}: 60 su 60 passaggi \w{0,940; 1} per geometria, radar a 360°, da 42\,950 a 44\,051 campioni per geometria, anche con frammenti più veloci del dichiarato (\textit{dim} nel modello: resistenza dell'aria, vento nullo).
\item \Sv\ \textbf{Radar accecato fino a 2 s insieme alle camere} (modello pessimista, \textit{ip}): mai «libero» \w{0; 0,060}, e il drone è rimasto nella regione in 59–60 passaggi su 60 \w{0,911; 0,997 nel caso peggiore}.
\item \Lv\ Con il radar accecato l'avviso al portatore prima di un frammento entro 10 m crolla (davanti 27–28 su 52, caso totale 0 su 59 \w{0; 0,061}; con la correzione in prova 59 su 59, vedi i limiti); col radar non disturbato arriva prima in tutti i casi simulati \Iv. Con un radar solo frontale un urto di lato o dietro non è visto (MoVi dice «non so»), e l'identità in virata scende a 43 su 60 \w{0,592; 0,815}. Il volume è prudente: l'avviso scatta anche senza frammenti vicini. Con l'urto a 10–30 m davanti il margine dell'avviso è sotto 0,25 s in 7 casi su 20. Il ritardo dell'evento non è presentabile. Il tracciatore usa in due punti la velocità vera del portatore. Nessun falso evento in 200 voli di controllo simulati senza clutter a grappoli \w{0; 0,019}.
\end{itemize}

\pagebreak
% ---------------------------------------------------------------- 5
\sez{5 · La collisione veloce: limite dichiarato, miglioria in prova}
\Lv\ La collisione veloce non è ancora dichiarata come evento. Una regola di contatto in prova (tracce risolte e mature, convergenza ≥ 10 m/s, fusione o sparizione, tracce nuove o Doppler fuori norma) la riconosce entro 2 s in 57 passaggi su 60 \w{0,863; 0,983} col radar a scansione elettronica, ma dà ancora falsi eventi in un quasi incrocio a 39 m, 2 su 60 \w{0,009; 0,114}, e un evento prima dell'urto; col radar rotante scende a 34 su 60 \w{0,441; 0,684}. \textbf{Non è adottata.} Si adotta solo con almeno 54 su 60 e zero falsi eventi su semi nuovi e almeno 10 geometrie. Per i droni lenti (8 m/s), con l'integratore stabile, l'evento arriva entro 2 s in 48 su 60 \w{0,682; 0,882}, con 3 eventi aperti prima dell'urto \Lv; i frammenti restano nel volume in 60 su 60 \w{0,940; 1} \Sv.

% ---------------------------------------------------------------- 6
\sez{6 · Sensori allo stato dell'arte e calcolo dedicato: stime, non misure}
{\fontsize{7.9}{9.4}\selectfont
\begin{tabularx}{\linewidth}{@{}>{\bfseries\RaggedRight}p{3.6cm}>{\RaggedRight}p{6.2cm}L@{}}
\toprule
Componente & Che cosa assumiamo (fonte) & Che cosa cambia \\
\midrule
Radar a scansione elettronica, banda Ku & classe Echodyne EchoShield: tracce a 10 Hz (stampa, 2026), campo 120×80° (produttore); portate per bersaglio non verificate & la regione dell'aereo veloce non si perde (§\,2) \\
Radar rotante & un giro al secondo (1 Hz) & la continuità dell'aereo veloce si perde, ed è dichiarata \\
Camere visibile e termica & termica «oltre 25 km» su bersaglio grande (dato del produttore) & accecate da fiammata e fumo: «non affidabili» \\
ADS-B, Remote ID & ≥ 1 messaggio/s, latenza ≤ 2,0 s (14 CFR 91.227) & solo chi trasmette \\
Calcolo dedicato (FPGA o ASIC) & CFAR su FPGA in 0,6 µs; distanza-Doppler 85× più rapida di un ARM; totale stimato \textbf{0,08–1,2 ms} per aggiornamento (\textit{ip}, per analogia) & la leva è la rivisita del sensore: 10 Hz → 0,62 s, 40 Hz → 0,17 s al primo avviso (\textit{calc} da \textit{ip}) \\
\bottomrule
\end{tabularx}}
\nota{Analogia con i chip per modelli linguistici: circa 17\,000 token/s contro circa 230 su una GPU H200 (dato dell'azienda, prove interne), valida solo per i carichi limitati dalla memoria. Per il passo di Kalman un FPGA è più lento in latenza (725 contro 250 ns) e circa 4,5 volte più rapido in portata (CERN). Verso un operatore umano la leva è la direzione dell'avviso: 5,4 s invece di 12,5 s (FAA AC 90-48E), 2005 m invece di 3312 m a 250 m/s (\textit{calc}).}

% ---------------------------------------------------------------- 7
\sez{7 · Lo stesso nucleo altrove: il sito con il rullaggio e il corto raggio}
\begin{itemize}
\item \Sv\ \textbf{Sito aerospaziale sorvegliato.} Un drone a 15–21 m/s dietro un edificio è annunciato in tempo in 200 su 200 \w{0,981; 1}; mai «libero» sul drone con la mappa dichiarata, 0 su 171 \w{0; 0,022}. L'avviso si spegne solo guardando tutta la zona: 0 su 600 \w{0; 0,006}. Il drone che rallenta o si ferma resta nell'avviso in 188 su 200 \w{0,898; 0,965} \Iv.
\item \Sv\ \textbf{Aereo in rullaggio dietro un hangar:} con tre sensori fusi l'avviso arriva in tempo in 200 su 200 \w{0,981; 1}; \Lv\ col solo radar in 50 su 200 \w{0,195; 0,314}. \Lv\ Per gli aerei in volo l'avviso al sito non è tempestivo: servono 3410 m (\textit{calc}).
\item \Sv\ \textbf{Corto raggio}, camera di profondità fino a 3 m: una persona nascosta non è mai rimasta senza avviso, 0 fotogrammi su 4\,324 \w{0; 0,001}. Per una persona non vedente è un aiuto d'avviso, mai provato con persone: non sostituisce il bastone né il cane guida.
\end{itemize}

% ---------------------------------------------------------------- 8
\sez{8 · I requisiti della traccia}
{\fontsize{7.9}{9.4}\selectfont
\begin{tabularx}{\linewidth}{@{}>{\bfseries\RaggedRight}p{4.1cm}L@{}}
\toprule
La traccia chiede & Che cosa fa MoVi, e con quale forza \\
\midrule
Continuità durante le interruzioni & regione mantenuta dopo urto, cecità e collisione vicina: 59–60 su 60 per cella col radar a scansione elettronica o a 360° \Sv; col radar rotante a 1 Hz 39 su 60 \Lv; confidenza nel tempo in forma chiusa, tarata su traiettorie pubbliche con sensore simulato \Iv, in virata \Lv \\
Distinguere osservato e incerto & tre risposte; «libero» solo se osservato \Sv; registro «osservato» o «dedotto» \\
Perdita di affidabilità & sensori saturi o accecati dichiarati «non affidabili» \Sv\ (innesco ottimista) \\
Eventi improvvisi & avaria dedotta entro 1 s \Iv; collisione veloce non dichiarata \Lv; volume dei detriti \Sv \\
Feedback tempestivi, latenza & avviso di prossimità ≥ 10 s prima \Iv; calcolo dedicato stimato, non misurato \\
Classificazione su più classi & supposta esistente, oggi etichette a mano \Nv; nessuna etichetta riduce un avviso \\
\bottomrule
\end{tabularx}}

% ---------------------------------------------------------------- 9
\sez{9 · Da proposta a sistema attuato: che cosa manca, e come lo si prova}
{\fontsize{7.9}{9.4}\selectfont
\begin{tabularx}{\linewidth}{@{}>{\RaggedRight}p{4.6cm}>{\RaggedRight}p{6.4cm}L@{}}
\toprule
Limite & Miglioria (che cosa e come) & Prova che la decide \\
\midrule
Collisione veloce non dichiarata & Doppler sulla cella unita, riagganci esclusi, radar vicino & ≥ 54/60 con zero falsi eventi, ≥ 10 geometrie \\
Confidenza in virata 0,73 & ricalibrazione condizionata a una manovra rilevata in modo causale; divisione per rotta & virata ≥ 0,85 (smentita se < 0,80) \\
Radar accecato: avviso al portatore & avviso «evento vicino non localizzato», in prova: 59/59, falsi 11/200; poi localizzarlo (direzione del lampo, Doppler) & avviso ≥ 0,9 con falsi ≤ 0,05 \\
Volume dei detriti prudente & velocità di ogni frammento dal Doppler; grappoli tracciati & avvisi falsi ≤ 10\,\% con contenimento al 100\,\% \\
Velocità vera del portatore nel tracciatore & velocità dalla navigazione inerziale & stesse soglie di §\,4 \\
Innesco della fiammata ottimista & saturazione dall'istogramma della camera & ≤ 1 ciclo su registrazioni vere \\
Una geometria per scenario; clutter a grappoli & ≥ 10 geometrie; stormi, pioggia, chaff simulati & stesse soglie \\
Calcolo, sensori, classificatore & FPGA reale; sensore fisico; classificatore per classe; prove sul campo & ≤ 1,2 ms per aggiornamento; prove sul campo \\
\bottomrule
\end{tabularx}}

\vspace{1pt}
\nota{\textbf{Dove sono le prove.} Formule e dimostrazioni nel modello, schede dei sensori, registri delle prove, intervalli e bibliografia nel documento tecnico. Il prototipo simulato produce i registri e il filmato con gli stessi comandi; il controllo delle regole sui registri usa lo stesso codice del nucleo (non indipendente).}

\end{document}
